% بخش: روش‌ها
% فصل: اثبات‌ها
% بخش فرعی: آغازکردن اثبات

\documentclass[../../../include/open-logic-section]{subfiles}

\begin{document}

\olfileid{mth}{prf}{str}

\olsection{آغاز یک اثبات}

اما اصلاً از کجا باید آغاز کرد؟

چیزی به شما داده شده است تا اثباتش کنید؛ پس باید آخرین چیزی باشد که
در اثبات ذکر می‌شود (البته می‌توانید در آغاز \emph{اعلام کنید} که
قصد اثباتش را دارید، اما نباید آن را همچون فرض به کار ببرید).
آنچه را می‌کوشید اثبات کنید پایینِ برگه‌ای تازه بنویسید؛ این‌گونه
هدفتان را از نظر دور نمی‌کنید.

سپس شاید فرض‌هایی داشته باشید که بتوانید به کار برید (وقتی در بخش
بعد دربارهٔ \emph{نوع} اثباتتان سخن بگوییم، این مطلب روشن‌تر
می‌شود). آن‌ها را بالای صفحه بنویسید و حتماً مشخص کنید که فرض‌اند
(یعنی اگر $p$ را فرض می‌کنید، بنویسید «فرض کنید $p$»). سرانجام،
شاید پرسش شامل تعریف‌هایی باشد که باید بدانید. ممکن است از شما
خواسته شود تعریف مشخصی را به کار ببرید، یا فرض‌ها و نتیجه‌ای که
در پی آن هستید تعریف‌های گوناگونی داشته باشند.
\emph{آن‌ها را بنویسید و مطمئن شوید معنایشان را می‌فهمید.}

شیوهٔ سامان‌دادن اثباتتان به صورت پرسش نیز بستگی دارد. بخش بعد
جزئیات سامان‌دادن اثبات را بر پایهٔ نوع جمله ارائه می‌کند.

\end{document}
